% PDF p. 256
\AppendixSection{5A}{测地三角形的高斯映射与球面几何}
\textbf{思路}\quad 高斯映射把曲率积分转化为球面三角形的面积。在单位球面 $S^2$ 上，测地三角形（边为大圆弧）的面积等于其球面角盈（spherical excess）$(\theta_1+\theta_2-\theta_3)$，这是吉拉尔定理（Girard's theorem）。整理各项便得证明。\translatornote{原书在此写出角盈 $(\theta_1+\theta_2-\theta_3)$，下一式还将内角和加到积分上。按本附录后面的式~\eqref{eq:5A8} 和式~\eqref{eq:5A10}，角盈是 $\theta_1+\theta_2+\theta_3-\pi$。}

\textbf{第 1 步}\quad 设 $T$ 是定向曲面 $(M,g)$ 上的一个测地三角形，内角为 $(\theta_1,\theta_2,\theta_3)$。证明
\[
  \iint_T K\,dA+(\theta_1+\theta_2+\theta_3)=\pi.
\]
先证明：对于定向曲面 $M\subset\R^3$ 的任一正则曲面片 $T$，若高斯映射
\[
  N\colon M\to S^2,\qquad
  N(p)=\left.\frac{X_u\times X_v}{\lVert X_u\times X_v\rVert}\right|_p
\]
为单射，即 $N$ 在 $T$ 上没有自身重叠，则
\[
  \operatorname{Area}(N(T))=\iint_T\lvert\det(dN_p)\rvert\,dA.
\]

% PDF p. 257
下面用换元证明上面的等式。直观上，关键在于相对于曲面的面积元素，法向量映射 $N$ 的雅可比行列式为 $\lvert\det(dN_p)\rvert$。
\begin{itemize}
  \item \textbf{设定：一张局部参数化}

  设
  \[
    X(u,v)=\bigl(x(u,v),y(u,v),z(u,v)\bigr),\qquad
    (u,v)\in\Omega\subset\R^2,
  \]
  为正则坐标卡，其像恰为曲面片 $T=X(\Omega)$。写出一阶偏导 $X_u,X_v\in\R^3$，令
  $W=\lVert X_u\times X_v\rVert>0$。$T$ 上的诱导面积元素为
  \begin{equation*}
    dA=W\,du\,dv.
    \tag{5A1}\label{eq:5A1}
  \end{equation*}
  \item \textbf{坐标下的高斯映射}

  将 $N$ 与 $X$ 复合，得到
  \[
    N(u,v)=\frac{X_u\times X_v}{W},\qquad N\colon\Omega\to S^2.
  \]
  计算一阶偏导 $N_u=\partial_uN$，$N_v=\partial_vN$。
  \item \textbf{高斯像的面积}

  在球面上，参数曲面片 $N(\Omega)$ 的面积元素为
  \[
    dA_{S^2}=\lVert N_u\times N_v\rVert\,du\,dv.
  \]
  因此
  \begin{equation*}
    \operatorname{Area}(N(T))=\iint_\Omega \lVert N_u\times N_v\rVert\,du\,dv.
    \tag{5A2}\label{eq:5A2}
  \end{equation*}
  我们需要 $\lVert N_u\times N_v\rVert$ 的显式表达式。下面推导该式。
  \item \textbf{把 $\lVert N_u\times N_v\rVert$ 与 $\det(dN_p)$ 联系起来}

  固定 $(u,v)$。微分 $dN_p\colon T_pM\to T_{N(p)}S^2$ 在切基 $\{X_u,X_v\}$ 上的作用为
  \[
    dN_p(X_u)=N_u,\qquad dN_p(X_v)=N_v.
  \]
  对任意由列向量 $Ae_1,Ae_2$ 表示的线性映射 $A\colon\R^2\to\R^2$，欧氏面积的放大因子为
  \[
    \frac{Ae_1\times Ae_2}{e_1\times e_2}.
  \]
  定义
  \[
    e_1=X_u,\quad e_2=X_v,\quad Ae_1=dN_p(X_u)=N_u,\quad
      Ae_2=dN_p(X_v)=N_v.
  \]
  于是
  \[
    \frac{\lVert Ae_1\times Ae_2\rVert}{\lVert e_1\times e_2\rVert}
       =\frac{\lVert N_u\times N_v\rVert}{\lVert X_u\times X_v\rVert}.
  \]
% PDF p. 258
  由换元定理，$dN_p$ 的雅可比行列式（绝对值）满足
  \begin{equation*}
    \begin{gathered}
      \lvert\det(dN_p)\rvert
         =\frac{\lVert N_u\times N_v\rVert}{\lVert X_u\times X_v\rVert},\\
      \text{从而}\quad\lVert N_u\times N_v\rVert
         =\lvert\det(dN_p)\rvert\lVert X_u\times X_v\rVert
         =\lvert\det(dN_p)\rvert W.
    \end{gathered}
    \tag{5A3}\label{eq:5A3}
  \end{equation*}
  \item \textbf{代入面积积分}

  将式~\eqref{eq:5A1} 和式~\eqref{eq:5A3} 代入式~\eqref{eq:5A2}，得
  \begin{equation*}
    \operatorname{Area}(N(T))
      =\iint_\Omega\lvert\det(dN_p)\rvert W\,du\,dv
      =\iint_\Omega\lvert\det(dN_p)\rvert\,dA.
    \tag{5A4}\label{eq:5A4}
  \end{equation*}
  由于 $X(\Omega)=T$，式~\eqref{eq:5A4} 右端的积分给出
  \begin{equation*}
    \operatorname{Area}(N(T))=\iint_T\lvert\det(dN_p)\rvert\,dA.
    \tag{5A5}\label{eq:5A5}
  \end{equation*}
  原书援引式 (5.50) 得\translatornote{此处应参见第 5.1.4.2 节对形状算子的定义和式~\eqref{eq:5.20}。原书所引的式 (5.50) 是雅可比方程。}
  \begin{equation*}
    S=-dN,\qquad K=\det S.
    \tag{5A6}\label{eq:5A6}
  \end{equation*}
  将式~\eqref{eq:5A6} 代入式~\eqref{eq:5A5}，得
  \begin{equation*}
    \operatorname{Area}(N(T))=\iint_T K\,dA,
    \tag{5A7}\label{eq:5A7}
  \end{equation*}
  其符号与所选的定向相合。\translatornote{式~\eqref{eq:5A5} 的普通面积对应 $\iint_T|K|\,dA$。若要写成 $\iint_TK\,dA$，应采用带定向的面积，或另有 $K\ge0$ 的条件。}

  \textbf{吉拉尔恒等式}\par
  \textbf{吉拉尔定理}\quad 一个球面三角形的内角为 $\theta_1,\theta_2,\theta_3$ 时，面积是
  \begin{equation*}
    \theta_1+\theta_2+\theta_3-\pi.
    \tag{5A8}\label{eq:5A8}
  \end{equation*}
  应用吉拉尔定理，得到
  \begin{equation*}
    \operatorname{Area}(N(T))=\theta_1+\theta_2+\theta_3-\pi.
    \tag{5A9}\label{eq:5A9}
  \end{equation*}
  合并式~\eqref{eq:5A7} 与式~\eqref{eq:5A9}，得到
  \begin{equation*}
    \iint_TK\,dA+\pi=\theta_1+\theta_2+\theta_3,
    \tag{5A10}\label{eq:5A10}
  \end{equation*}
% PDF p. 259
  等价地，
  \begin{equation*}
    \iint_TK\,dA+\varnothing_1+\varnothing_2+\varnothing_3=2\pi,
    \tag{5A11}\label{eq:5A11}
  \end{equation*}
  其中 $\varnothing_i=\pi-\theta_i$ 为外角。
  \item \textbf{测地线在高斯映射下的像}

  设 $\gamma(s)$ 是 $T$ 的一条单位速率测地边。因为 $\ddot\gamma$ 是 $M$ 的法向量，法向量场 $N(\gamma(s))$ 始终位于过原点且与 $\dot\gamma(0)$ 垂直的固定平面内。这个平面与单位球面的交集是大圆。因此，$T$ 的每条测地边在 $N$ 下都映为一段大圆弧。\translatornote{测地线条件只保证 $N(\gamma(s))$ 与当前的 $\dot\gamma(s)$ 垂直，不能推出与固定向量 $\dot\gamma(0)$ 垂直。下文抛物面子午线是特殊例子，不能据此断定一般高斯像 $N(T)$ 是球面测地三角形。}

  在抛物面片 $z=\frac12(x^2+y^2)$ 上，所选测地边是一条子午线，因而确是测地线。其高斯映射为
  \[
    N(p)=\frac{(-x,-y,1)}{\sqrt{x^2+y^2+1}}.
  \]
  考虑竖直平面 $y=0$。这个平面与单位球面（浅青色）相交成一个大圆。因此，高斯映射将该测地边映为大圆的一段弧，用来说明“$T$ 的每条测地边在 $N$ 下都映为大圆弧”。由此，原书得出 $N(T)$ 是以大圆弧为边的球面三角形。
  \item \textbf{曲面上的角与球面上的角}

  考虑在顶点 $p$ 相交的两条测地边，其内角为 $\theta$。设 $v,w$ 为向外的切向量。由于 $N$ 垂直于 $v,w$，有
  \[
    \langle v,w\rangle=\cos\theta,\qquad
    \langle N_v,N_w\rangle=\cos\varnothing,
  \]
  其中 $\varnothing$ 为球面三角形在 $N(p)$ 处的内角。\translatornote{计算角的余弦需要先归一化向量。一般的 $dN_p$ 不保角，本段的两项公式不能推出两个角相等，故不能用同一组内角得到式~\eqref{eq:5A9}。}
\end{itemize}
